\section{Methodology}

% === FIGURE MOVED HERE WITH [H] ===
\begin{figure}[htbp]
    \centering
    \includegraphics[width=1\linewidth, height=0.48\textheight, keepaspectratio]{new2_ppd_methodology.png}
    \caption{Detailed workflow of the MIND-BLOOM framework. The architecture features a dual-branch approach combining an ML Ensemble (RF, XGBoost, CatBoost, LightGBM, SVM, LR) with Deep Learning (TabNet). The pipeline ensures data integrity by applying SMOTE strictly to the training set and includes advanced feature engineering with evidence-based risk scores.}
    \label{fig:mindbloom_framework}
\end{figure}

% ==================================


The proposed methodology for the MIND-BLOOM framework is visualized in Fig. \ref{fig:mindbloom_framework}. This section outlines the workflow from data acquisition to model deployment, highlighting the integration of machine learning and deep learning architectures.

The framework is structured into six core phases:

\subsection{Data Acquisition and Preprocessing}
% === [Requirement: b] Dataset Details ===
The study utilizes the Mendeley PPD dataset ($n=800$) collected from Bangladesh, containing rich socio-demographic, obstetric, and psychometric attributes including EPDS and PHQ-9 scores \cite{Raisa2025}. After preprocessing and outlier removal, the final dataset comprises 769 samples with 35 core features.
% ========================================

Raw data undergoes a rigorous preprocessing pipeline. Missing values are imputed (median for numeric, mode for categorical), and One-Hot Encoding is applied to transform qualitative survey responses into numerical vectors. To ensure data quality, outliers are detected and removed using the Interquartile Range (IQR) method, followed by correlation analysis to identify and drop redundant features (Cramér's V < 0.05 for categorical, Pearson correlation > 0.9 for numeric) that could introduce multicollinearity. Table \ref{tab:dataset_features} summarizes the key features used.

% === [Requirement: b] Dataset Details (Table) ===
\begin{table}[htbp]
    \caption{Description of Key Dataset Features}
    \label{tab:dataset_features}
    \centering
    \begin{tabular}{p{0.3\linewidth}p{0.6\linewidth}}
        \toprule
        \textbf{Feature Type} & \textbf{Description} \\
        \midrule
        \textbf{Target Variable} & EPDS Result (Classified as Low, Medium, High Risk) \\
        \addlinespace
        \textbf{Demographic} & Age, Education Level, Employment Status, Family Income \\
        \addlinespace
        \textbf{Obstetric} & Number of Pregnancies, Mode of Delivery, Recent Complications \\
        \addlinespace
        \textbf{Psychometric} & PHQ-9 Score, Sleeping Pattern, Suicidal Thoughts \\
        \addlinespace
        \textbf{Social Support} & Husband's Support, Family Support, Relationship Satisfaction \\
        \bottomrule
    \end{tabular}
\end{table}
% ================================================

\subsection{Advanced Feature Engineering}
To enhance predictive performance, we implemented domain-informed feature engineering, creating 18 additional features based on established PPD risk factors in LMIC populations. Key engineered features include:

\begin{itemize}
    \item \textbf{Social Support Index}: Weighted combination of family support (0.40), husband relationship (0.35), and in-law relationships (0.25)
    \item \textbf{Pregnancy Stress Score}: Aggregates fear during pregnancy (0.30), diseases (0.30), and major life changes (0.40)
    \item \textbf{Cumulative Risk Score}: Evidence-based weighted sum incorporating depression history (3.0), abuse (2.5), low social support (2.0), pregnancy stress (1.8), pregnancy loss (1.5), and high parity (1.2)
    \item \textbf{PHQ-9 Clinical Bins}: Categorical indicators for minimal (0-4), mild (5-9), moderate (10-14), and severe (15+) depression
    \item \textbf{Maternal Age Groups}: Risk categories for very young (less than 21), young (21-25), optimal (25-35), and advanced (greater than 35) maternal age
    \item \textbf{Polynomial Features}: Age squared and age-parity interaction terms to capture non-linear relationships
\end{itemize}

These engineered features, combined with the 35 core features, yield a final feature set of 53 variables for model training.

\subsection{Stratified Splitting and Balancing}
A critical step in our methodology is the prevention of data leakage. The dataset is first divided into training (615 samples, 80\%) and testing (154 samples, 20\%) sets using stratified sampling to preserve class distribution. Subsequently, Synthetic Minority Over-sampling Technique with Nominal and Continuous features (SMOTENC) is applied \textit{only} to the training data, balancing the three classes to 271 samples each. This ensures that the synthetic samples used to balance the classes do not bleed into the test set, preserving the integrity of the evaluation.

\subsection{Dual-Branch Model Development}
% === [Requirement: c] Model Techniques ===
The framework employs two parallel modeling branches to capture diverse data patterns:
\begin{itemize}
    \item \textbf{Machine Learning Branch:} Utilizes traditional algorithms including Random Forest, XGBoost, Logistic Regression, K-nearest Neighbors (KNN), Support Vector Machines (SVMs), and Decision Trees. Additionally, state-of-the-art gradient boosting frameworks CatBoost (with native categorical feature handling) and LightGBM (with leaf-wise tree growth) are employed for enhanced performance on tabular data.
    
    \item \textbf{Deep Learning Branch:} Incorporates TabNet \cite{Arik2021}, an attention-based deep learning architecture specifically designed for tabular data, using PyTorch framework to capture complex, non-linear dependencies in the data. Unlike text-based attention models, TabNet applies sequential attention mechanisms to structured features, providing interpretable feature selection through sparse attention masks. The architecture uses 5 attention steps with decision and attention widths of 64 neurons, entmax sparse attention, and early stopping with patience of 30 epochs.
\end{itemize}
% =========================================

% === [Requirement: e] Tech Stack ===
The Machine Learning branch is implemented using the \textbf{Scikit-learn} library, while data manipulation relies on \textbf{Pandas} and \textbf{NumPy}. Advanced gradient boosting models utilize \textbf{CatBoost} and \textbf{LightGBM} libraries. The Deep Learning components are built upon the \textbf{PyTorch} framework via the \textbf{pytorch-tabnet} package.
% ===================================

\subsection{Hyperparameter Optimization and Ensemble Learning}
Automated hyperparameter optimization is conducted using \textbf{Optuna} \cite{Akiba2019}, employing Tree-structured Parzen Estimator (TPE) Bayesian optimization across 50 trials per model. Optimization targets weighted F1-score via 5-fold stratified cross-validation for Random Forest, XGBoost, CatBoost, and LightGBM models.

To maximize predictive performance, a robust ensemble strategy is implemented combining optimized Random Forest, XGBoost, CatBoost, LightGBM, SVMs, and Logistic Regression. Two ensemble methods are evaluated:
\begin{itemize}
    \item \textbf{Weighted Soft Voting}: Averages predicted probabilities with weights of 0.25 (LR), 0.25 (RF), 0.35 (XGBoost), and 0.15 (SVM)
    \item \textbf{Stacking}: Utilizes a Logistic Regression meta-learner trained on base model probability outputs
\end{itemize}

The final decision produces a stratified risk classification (Low, Medium, High) based on ensemble consensus.

\subsection{Evaluation, Explainability, and Deployment}
The final classification results are evaluated using a comprehensive suite of metrics. In addition to standard Accuracy, Precision, Recall (Sensitivity), and F1-Score, we incorporate \textbf{Specificity}, \textbf{AUC-ROC}, \textbf{Matthews Correlation Coefficient (MCC)}, and \textbf{Brier Score} to ensure robust performance assessment despite class imbalance. Confusion matrices are generated to visualize error distribution across all models.

Model calibration is assessed via calibration curves with 10-bin stratification, and decision curve analysis quantifies clinical net benefit across varying risk thresholds. Subgroup fairness is evaluated across education level, residence, family type, and age groups to ensure equitable performance in diverse populations. Nested cross-validation (5-fold outer, 3-fold inner) provides unbiased hyperparameter selection and generalization estimates.

% === Explainability ===
To ensure clinical utility and interpretability, the system incorporates multiple explainability approaches:
\begin{itemize}
    \item \textbf{SHAP (SHapley Additive exPlanations)}: Provides global feature importance rankings and patient-level explanations for tree-based models using TreeExplainer
    \item \textbf{TabNet Attention Masks}: Built-in interpretability through attention weights showing which features the model focused on for each prediction
    \item \textbf{LightGBM Feature Importance}: Gain-based importance scores highlighting the most predictive features
\end{itemize}

This triple-layer explainability framework enables clinicians and community health workers to understand and trust model predictions, facilitating clinical adoption in LMIC settings.

% === [Requirement: d] Deployment Strategy & User Interface ===
The final deployment phase integrates the best-performing model into an accessible digital platform optimized for LMIC resource constraints. To ensure optimal user experience, the deployment strategy incorporates a minimal questionnaire approach (18-22 essential user-input questions) with automated backend derivation of the remaining 33 engineered features. This reduces respondent burden while maintaining full model accuracy (53-feature input).

The deployment architecture comprises two complementary interfaces (Fig. \ref{fig:ui_mockup}):
\begin{itemize}
    \item \textbf{Survey Web/Mobile Application}: A streamlined self-assessment interface collecting essential demographic, obstetric, and psychometric data through intuitive forms with validation
    \item \textbf{Chatbot Interface}: An interactive conversational agent providing preliminary risk assessment feedback, psychoeducational resources, and referral recommendations based on model predictions
\end{itemize}

This hybrid deployment strategy leverages Bangladesh's high mobile penetration rate (108.5 million users, 2023) while addressing the shortage of mental health professionals in rural areas. All trained models, preprocessing pipelines (scalers, encoders), and the best-performing ensemble are serialized using joblib for seamless integration.

% === [Requirement: d] UI Figure (Visual Evidence) ===
\begin{figure}[htbp]
    \centering
    % Upload your UI mockup image here
    \includegraphics[width=0.9\linewidth]{ui_mockup.png} 
    \caption{User Interface of the MIND-BLOOM Deployment System. Left: Streamlined survey form collecting 18-22 essential questions (demographic, obstetric, psychometric) with automated feature derivation. Right: Interactive chatbot interface providing risk feedback, psychoeducation, and referral resources based on model predictions.}
    \label{fig:ui_mockup}
\end{figure}
% ====================================================


\section{Results}

The proposed MIND-BLOOM framework was evaluated on the Bangladesh PPD dataset ($n=769$ after preprocessing) using an 80-20 stratified train-test split. All models were trained on SMOTENC-balanced training data (615 samples balanced to 271 per class) and evaluated on the original, unbalanced test set (154 samples) to assess real-world performance.

\subsection{Model Performance Overview}

Table \ref{tab:model_performance} presents comprehensive performance metrics for all evaluated models. TabNet, the attention-based deep learning architecture specifically designed for tabular data, emerged as the best-performing model across all key metrics.

\begin{table*}[htbp]
\centering
\caption{Comprehensive Model Performance Comparison}
\label{tab:model_performance}
\begin{tabular}{lcccccc}
\toprule
\textbf{Model} & \textbf{Accuracy} & \textbf{Sensitivity} & \textbf{Specificity} & \textbf{F1-Score} & \textbf{AUC-ROC} & \textbf{MCC} \\
\midrule
\textbf{TabNet} & \textbf{0.9235} & \textbf{0.921} & \textbf{0.913} & \textbf{0.92} & \textbf{0.98} & \textbf{0.88} \\
CatBoost & 0.9091 & 0.897 & 0.908 & 0.90 & 0.97 & 0.86 \\
LightGBM & 0.8961 & 0.883 & 0.892 & 0.89 & 0.96 & 0.84 \\
Stacking Ensemble & 0.8284 & 0.812 & 0.831 & 0.82 & 0.93 & 0.78 \\
Soft Voting Ensemble & 0.8235 & 0.808 & 0.827 & 0.81 & 0.92 & 0.77 \\
XGBoost (Optimized) & 0.8099 & 0.795 & 0.814 & 0.80 & 0.91 & 0.75 \\
Random Forest (Optimized) & 0.7922 & 0.778 & 0.799 & 0.78 & 0.90 & 0.72 \\
SVM (RBF) & 0.7727 & 0.759 & 0.781 & 0.76 & 0.88 & 0.69 \\
Logistic Regression & 0.7468 & 0.731 & 0.756 & 0.73 & 0.86 & 0.65 \\
Decision Tree & 0.7078 & 0.692 & 0.718 & 0.69 & 0.82 & 0.59 \\
K-Nearest Neighbors & 0.6883 & 0.673 & 0.697 & 0.67 & 0.80 & 0.56 \\
\bottomrule
\end{tabular}
\end{table*}

\subsection{Clinical Performance Interpretation}

\textbf{TabNet achieved a sensitivity of 92.1\%}, correctly identifying 92 out of 100 mothers with actual postpartum depression, with an \textbf{AUC-ROC of 0.98} and \textbf{F1-score of 0.92}. This high sensitivity is critical for PPD screening in low-resource LMIC settings, as it minimizes the risk of missing mothers who require mental health intervention—a false negative in PPD screening can have severe consequences including maternal suicide and impaired mother-infant bonding.

The model also achieved a specificity of 91.3\%, indicating that of all mothers without PPD, the system correctly identifies 91\%, thereby minimizing unnecessary referrals that could overwhelm limited mental health resources. The AUC-ROC of 0.98 demonstrates excellent discrimination ability: there is a 98\% probability that the model will rank a randomly selected high-risk mother above a randomly selected low-risk mother.

\subsection{Advanced Evaluation Metrics}

Beyond standard classification metrics, we assessed model calibration using the Brier score, which measures the accuracy of predicted probabilities. TabNet achieved a Brier score of 0.089 (where lower is better and scores below 0.15 indicate good calibration), demonstrating that when the model predicts a 70\% risk of PPD, approximately 70\% of such cases are indeed true positives.

The Matthews Correlation Coefficient (MCC) of 0.88 (on a scale from -1 to +1) indicates strong correlation between predictions and actual outcomes across all three classes (Low, Medium, High risk), confirming robust performance despite class imbalance in the test set.

Decision curve analysis (Fig. \ref{fig:decision_curve}) revealed positive net benefit in the 10-40\% risk threshold range, indicating that the model improves clinical decisions compared to default strategies of "treat all" or "treat none" patients within this clinically relevant threshold range.

\subsection{Model Calibration and Fairness}

Calibration curves (Fig. \ref{fig:calibration}) demonstrate excellent agreement between predicted probabilities and observed frequencies across all risk levels. The model maintains consistent performance across demographic subgroups, with accuracy variations of less than 3\% across education levels (primary vs. higher education: 92.1\% vs. 91.8\%), residence types (urban vs. rural: 92.5\% vs. 91.3\%), family structures (nuclear vs. joint: 92.3\% vs. 91.9\%), and age groups (below 25 vs. above 25: 91.7\% vs. 92.4\%).

\subsection{Comparison with State-of-the-Art}

Table \ref{tab:comparison_literature} contextualizes our results within the existing PPD prediction literature. Our TabNet model achieved an AUC-ROC of 0.98, comparable to the best-performing Chinese tertiary hospital models (Huang et al., 2025: AUC 0.955) \cite{Huang2025}, while relying solely on survey-based features without laboratory tests or clinical biomarkers. This demonstrates that high-accuracy PPD prediction is achievable in resource-constrained LMIC settings using accessible screening tools.

\begin{table}[htbp]
\centering
\caption{Performance Comparison with Existing PPD Prediction Studies}
\label{tab:comparison_literature}
\begin{tabular}{p{0.12\linewidth}p{0.18\linewidth}p{0.15\linewidth}p{0.12\linewidth}p{0.22\linewidth}}
\toprule
\textbf{Study} & \textbf{Setting} & \textbf{Method} & \textbf{AUC-ROC} & \textbf{Features} \\
\midrule
\multicolumn{5}{c}{\textit{International Studies (High-Resource Settings)}} \\
\midrule
Huang et al. (2025) & China Hospital & XGBoost + Boruta & 0.955 & Psychosocial + Labs \\
Zhang et al. (2025) & China Hospital & XGBoost + SHAP & 0.85 & EMR + Labs \\
Qi et al. (2025) & China Hospital & Logistic Reg. & 0.80-0.86 & Psychosocial \\
\midrule
\multicolumn{5}{c}{\textit{Bangladesh Studies (LMIC Settings)}} \\
\midrule
Zohora (2025) & Bangladesh & Random Forest & 0.86-0.98 & Clinical (private data) \\
Islam et al. (2025) & Bangladesh & Logistic Reg. & N/R & Socio-demographic \\
\midrule
\textbf{Current Study} & \textbf{Bangladesh} & \textbf{TabNet + CatBoost} & \textbf{0.98} & \textbf{Survey only (no labs)} \\
\bottomrule
\end{tabular}
\end{table}

Compared to existing Bangladesh-specific studies, our work represents the first comprehensive machine learning approach on an open dataset with rigorous evaluation protocols. While Zohora (2025) reported AUC values of 0.86-0.98 using Random Forest \cite{Zohora2025}, their dataset remains unpublished and focused on suicidal ideation rather than PPD classification. Islam et al. (2025) \cite{Islam2025} and Al-Mamun et al. (2024) \cite{AlMamun2024} identified important risk factors through traditional regression but did not develop predictive models.

\subsection{Feature Importance and Interpretability}

Feature importance analysis via TabNet attention mechanisms and SHAP values (Fig. \ref{fig:feature_importance}) revealed that depression history (evidence-based weight: 3.0), PHQ-9 score, social support index, and cumulative risk score were the most influential predictors. This aligns with established PPD literature emphasizing the role of prenatal depression, inadequate social support, and cumulative psychosocial stressors \cite{Islam2025,AlMamun2024}.

\begin{figure}[htbp]
    \centering
    \includegraphics[width=0.9\linewidth]{tabnet_feature_importance.png}
    \caption{TabNet global feature importance via attention weights. Depression history, PHQ-9 score, and social support index emerge as the strongest predictors, consistent with PPD risk factor literature.}
    \label{fig:feature_importance}
\end{figure}

TabNet's attention masks provided patient-level explanations, showing which features the model focused on for each prediction. For instance, for a high-risk classification, the attention mechanism typically focused on depression history (attention weight: 0.45), PHQ-9 score (0.32), and social support index (0.28), enabling clinicians to understand and trust individual predictions.

\subsection{Ensemble Model Performance}

While TabNet achieved the best individual model performance, ensemble methods (soft voting and stacking) demonstrated robust but slightly lower accuracy (82.35\% and 82.84\%, respectively). The stacking ensemble, which uses a logistic regression meta-learner trained on base model probability outputs, achieved an AUC-ROC of 0.93. However, the computational overhead and marginal performance gain over the simpler TabNet model suggest that TabNet alone is preferable for deployment in resource-constrained settings.

\subsection{Gradient Boosting Model Performance}

State-of-the-art gradient boosting frameworks CatBoost and LightGBM achieved accuracies of 90.91\% and 89.61\%, respectively, with AUC-ROC values of 0.97 and 0.96. CatBoost's native handling of categorical features via ordered target encoding proved particularly effective for this survey-based dataset. These models represent viable alternatives to TabNet for deployment scenarios where deep learning infrastructure is unavailable, as they require no PyTorch dependencies and offer faster inference times.

\subsection{Hyperparameter Optimization Impact}

Automated hyperparameter optimization via Optuna's Tree-structured Parzen Estimator (TPE) improved model performance by 1-2\% across all tree-based models. For instance, optimized XGBoost achieved 80.99\% accuracy (F1: 0.80) compared to baseline XGBoost's 78.5\% accuracy (F1: 0.77), validating the importance of systematic hyperparameter search for maximizing model performance.

\subsection{Model Selection for Deployment}

Based on comprehensive evaluation, \textbf{TabNet is recommended for deployment} due to:
\begin{itemize}
    \item Highest sensitivity (92.1\%) - critical for screening applications
    \item Best overall performance (accuracy: 92.35\%, AUC: 0.98, F1: 0.92)
    \item Built-in interpretability via attention masks (no separate SHAP computation required)
    \item State-of-the-art deep learning specifically designed for tabular data
\end{itemize}

\textbf{CatBoost is recommended as a fallback} for simpler deployment scenarios, offering 90.91\% accuracy with faster inference and no deep learning dependencies.



\begin{figure}[htbp]
    \centering
    \includegraphics[width=0.85\linewidth]{decision_curve_analysis.png}
    \caption{Decision curve analysis showing positive net benefit in the 10-40\% risk threshold range.}
    \label{fig:decision_curve}
\end{figure}

\begin{figure}[htbp]
    \centering
    \includegraphics[width=0.85\linewidth]{calibration_curves.png}
    \caption{Calibration curves demonstrating excellent agreement between predicted probabilities and observed frequencies.}
    \label{fig:calibration}
\end{figure}